\documentclass[a4paper,10pt]{article}

\usepackage{tikz}
\usetikzlibrary{arrows.meta, positioning, calc}
\usepackage{geometry}
\usepackage{hyperref}
\usepackage{titling}

\geometry{margin=5mm}
\setlength{\droptitle}{-2.5em}

\hypersetup{
    colorlinks=true,
    urlcolor=blue
}

\title{Deterministic Failure in Display State Transition of Z f (C: Ver. 3.01)}
\author{
    ``CASE-A1: Schrödinger's Sisyphean Loop''
    \thanks{
        Etymology: Named after Schrödinger's cat to describe the non-deterministic
        state collapse triggered by shutter ``observation,'' and the myth of
        Sisyphus to reflect the futile loop of Info displays that prevents a
        return to Live View.
    }
}
\date{}

\begin{document}
\maketitle

\noindent\vspace{-2cm}

\noindent\hspace{-1.5cm}
\begin{tikzpicture}[
    scale=0.96,
    transform shape,
    node distance=1.5cm and 2.2cm,
    every node/.style={
        draw,
        rectangle,
        rounded corners,
        align=center,
        font=\scriptsize,
        minimum width=1.8cm,
        minimum height=0.6cm
    },
    normal/.style={->, >=Stealth, thick, green!50!black},
    abnormal/.style={->, >=Stealth, thick, green!50!black, dashed},
    critical/.style={->, >=Stealth, thick, green!50!black},
    label_style/.style={
        font=\tiny\itshape,
        fill=white,
        inner sep=0.5pt,
        outer sep=1pt
    },
    state_normal/.style={fill=blue!5},
    state_abnormal/.style={draw=red, fill=red!5, line width=1.2pt},
    state_liveview/.style={draw=blue, fill=blue!5, line width=1.2pt},
    state_menuo/.style={draw=blue, fill=yellow!16, line width=1.2pt},
    state_imenu/.style={draw=blue, fill=yellow!16},
    state_setup/.style={fill=gray!10, dashed}
]

% --------------------------------------------------
% State Placement
% --------------------------------------------------

\node[state_setup] (S0) at (-0.4,0) {S0\\LCD Off};
\node[state_liveview] (S1) at (-0.4,-2.5) {S1\\LiveView};
\node[state_normal] (S9) at (-0.4,-5.0) {S9\\Selfie};

\node[state_normal] (S2) at (-5.6,0) {S2\\Menu (D)};
\node[state_menuo] (S3) at (-5.6,-2.2) {S3\\Menu (O)};

\node[state_imenu] (S4) at (5.8,1.1) {S4\\iMenu (D)};
\node[state_imenu] (S5) at (5.8,-1.1) {S5\\Info (O)};

\node[state_abnormal] (S6) at (5.8,-7) {
    S6 Stuck in a loop\\
    Observed: LCD = Display 5 (Info)\\
    (Hypothesis: Live View routing inactive,\\
    display routed to INFO\_ONLY)
};

\node[state_abnormal] (S7) at (6.3,-9.3) {S7\\Info Hold};

\node[state_normal] (S10) at (5.5,-11.1) {S10\\Play (D)};
\node[state_normal] (S11) at (3.8,-1.1) {S11\\Play (O)};
\node[state_normal] (S12) at (7.8,0) {S12\\WB (D)};
\node[state_normal] (S13) at (7.8,-1.1) {S13\\WB (O)};

% --------------------------------------------------
% Normal Transitions
% --------------------------------------------------

\draw[normal, green!25!black!, line width=1.5pt]
    (S0) edge[bend left=50, pos=0.55]
    node[label_style, fill=black!10]
    {History-dependent\\on previous action/setting}
    (S1);

\draw[normal]
    (S1) edge[bend left=75, pos=0.65]
    node[label_style]{Close}
    (S0);

\draw[normal]
    (S1) edge[bend left=55]
    node[label_style]{180$^\circ$}
    (S9);

\draw[normal]
    (S9) edge[bend left=55]
    node[label_style]{Back}
    (S1);

\draw[normal]
    (S0) -- node[label_style, pos=0.6]{MENU}
    (S2);

\draw[normal]
    (S2) -- node[label_style]{Open}
    (S3);

\draw[normal, green!25!black!, line width=1.5pt]
    (S3) .. controls (-4.5,-4.6) and (-2.9,-3.2) ..
    node[label_style, pos=0.30, fill=black!9]{Half-press}
    (S1);

% --------------------------------------------------
% Persistent / Loop Transitions
% --------------------------------------------------

\draw[abnormal]
    (S5) -- node[label_style, pos=0.3]{Half-press shutter}
    (S6);

\draw[abnormal, green!25!black!, line width=1.5pt]
    (S3) .. controls (-10.3,-7) and (-7.6,-6.6) ..
    node[label_style, pos=0.085, fill=black!8]{Half-press}
    (S6);

\draw[abnormal]
    (S11) -- node[label_style, pos=0.45]{Playback}
    (S6);

\draw[abnormal]
    (S13) -- node[label_style, pos=0.45]{Fn Release}
    (S6);

\draw[abnormal]
    (S9) .. controls (0,-6) and (2,-6.6) ..
    node[label_style, pos=0.4]{Back}
    (S6);

\draw[abnormal]
    (S12) .. controls (9.8,0) and (9.8,-9.5) ..
    node[label_style, pos=0.47]{Fn Release}
    (S7);

\draw[abnormal]
    (S6) edge[bend left=45]
    node[label_style, pos=0.48]{Close}
    (S7);

\draw[abnormal]
    (S7) edge[bend left=55]
    node[label_style]{Open}
    (S6);

\draw[abnormal, dashed, line width=1.5pt]
    (S6) .. controls (1.9,-5.3) and (4.4,-0.5) ..
    node[label_style, pos=0.25]
    {Apparent reset\\(state not\\guaranteed\\clean)}
    (S0);

\draw[abnormal, dashed, green!25!black!, line width=1.5pt]
    (S0) .. controls (2.8,-0.7) and (0.8,-6.3) ..
    node[label_style, pos=0.42]
    {\\[0.05pt] Multiple steps\\(see transition table)\\
    \textit{e.g., Info displayed}\\
    \textit{on the LCD}}
    (S6);

\draw[normal]
    (S0) -- node[label_style, pos=0.55]{[i] button}
    (S4);

\draw[abnormal, dashed]
    (S4) .. controls (13,1.4) and (10,-10) ..
    node[label_style, pos=0.45]{[i] button}
    (S7);

% --------------------------------------------------
% Recovery Paths
% --------------------------------------------------

\draw[normal, line width=1.5pt, blue!90!black]
    (S6) .. controls (-9.2,-13.7) and (-4,-4) ..
    node[label_style, pos=0.256, sloped]
    {DISP button\\Disable Display 5\\Reset}
    (S1);

\draw[abnormal]
    (S7) .. controls (-1.0,-9.6) and (-3.0,-8) ..
    node[label_style, pos=0.11, sloped]
    {\\DISP button\\LCD Open}
    (S6);

\draw[critical, line width=1.5pt, purple!70!black]
    (S6) .. controls (-6.8,-10.8) and (-3.2,-7.0) ..
    node[label_style, pos=0.32, sloped]
    {DISP button\\(if assigned)}
    (S6);

% --------------------------------------------------
% Playback / White Balance Paths
% --------------------------------------------------

\draw[normal]
    (S7) -- node[label_style, pos=0.45]{Play}
    (S10);

\draw[normal]
    (S0) -- node[label_style, pos=0.7]{Fn}
    (S12);

\draw[abnormal]
    (S10) .. controls (-1.5,-10) and (-3.8,-7.5) ..
    node[label_style, pos=0.12]{Playback}
    (S6);

% --------------------------------------------------
% Diagram Annotations
% --------------------------------------------------

\node[
    label_style,
    font=\tiny,
    align=left,
    inner sep=3pt,
    fill=black!8,
    dashed
]
at (-5.75, -5.1) {
    \rule{0pt}{22pt}
    Step 5 / 28:\\
    Half-press shutter\\
    \textit{State collapse triggered}\\
    \textit{by observation}\\
    \textit{(Hidden Preconditions,}\\
    \textit{Interaction History)}
};

\node[
    draw=red,
    dashed,
    fill=yellow!5,
    text width=3.3cm,
    font=\tiny,
    align=left
]
at (-1.4,1.3) {
    \textbf{Note:}\\
    S0 does not guarantee\\
    a clean internal state.\\
    Hidden state may persist and\\
    influence future transitions.
};

\node[
    draw=red,
    dashed,
    fill=yellow!5,
    text width=8cm,
    font=\tiny,
    align=left,
    inner sep=3pt
]
at (-3.5,-11.0) {
    \textbf{Critical UI Loophole:}\\
    The DISP-related recovery behavior can be altered, reassigned,
    or effectively removed through Custom Settings without any explicit
    system warning. Under certain configurations, standard UI-defined
    recovery paths from persistent Info-display states may become
    unavailable through ordinary shooting operations.\\
    Additionally, power cycling or reinstalling the battery was not
    observed to reliably terminate these persistent states.
};

\node[
    draw=none,
    fill=black!3,
    text width=3.6cm,
    font=\tiny,
    align=left,
    inner sep=3pt
]
at (9,-11) {
    \textbf{Settings:}\\
    Monitor mode\\
    = Prioritize viewfinder (1 or 2)\\
    Automatic monitor display switch\\
    = On (when monitor docked)
};

% --------------------------------------------------
% Conceptual Illustration Elements
% --------------------------------------------------

\newcommand{\myfscale}{0.5}
\newcommand{\myfx}{1.2}
\newcommand{\myfy}{1.2}

\begin{scope}[shift={(\myfx,\myfy)}, scale=\myfscale]
    \draw (0,0) ellipse (1.5cm and 0.8cm);

    \begin{scope}[shift={(1.1,0.3)}, rotate=-10]
        \draw (0,0) circle (0.5cm);
        \draw (-0.3,0.4) -- (-0.4,0.8) -- (0,0.5);
        \draw (0.3,0.4) -- (0.5,0.7) -- (0.5,0.2);
        \draw (-0.3,0) arc (180:360:0.1);
        \draw (0.1,0) arc (180:360:0.1);
        \draw (0,-0.15) -- (0,-0.25);
        \draw (0,-0.25) arc (180:360:0.05);
        \draw (0,-0.25) arc (0:-180:0.05);
    \end{scope}

    \draw (-1.4,-0.3) .. controls (-2.5,-0.5) and (-2.5,1) .. (-1.5,0.6);
    \draw (-0.5,-0.8) arc (200:340:0.4);
\end{scope}

\newcommand{\mylscale}{0.26}
\newcommand{\mylx}{-5.2}
\newcommand{\myly}{-4.5}

\begin{scope}[shift={(\mylx,\myly)}, scale=\mylscale]
    \draw (0,0) ellipse (1.5cm and 0.8cm);

    \begin{scope}[shift={(1.1,0.3)}, rotate=-10]
        \draw (0,0) circle (0.5cm);
        \draw (-0.3,0.4) -- (-0.4,0.8) -- (0,0.5);
        \draw (0.3,0.4) -- (0.5,0.7) -- (0.5,0.2);
        \draw (-0.3,0) arc (180:360:0.1);
        \draw (0.1,0) arc (180:360:0.1);
        \draw (0,-0.15) -- (0,-0.25);
        \draw (0,-0.25) arc (180:360:0.05);
        \draw (0,-0.25) arc (0:-180:0.05);
    \end{scope}

    \draw (-1.4,-0.3) .. controls (-2.5,-0.5) and (-2.5,1) .. (-1.5,0.6);
    \draw (-0.5,-0.8) arc (200:340:0.4);
\end{scope}

\newcommand{\myrscale}{0.26}
\newcommand{\myrx}{-6.25}
\newcommand{\myry}{-4.5}

\begin{scope}[shift={(\myrx,\myry)}, scale=\myrscale, xscale=-1]
    \draw (0,0) ellipse (1.5cm and 0.8cm);

    \begin{scope}[shift={(1.1,0.3)}, rotate=-10]
        \draw (0,0) circle (0.5cm);
        \draw (-0.3,0.4) -- (-0.4,0.8) -- (0,0.5);
        \draw (0.3,0.4) -- (0.5,0.7) -- (0.5,0.2);
        \draw (-0.3,0) arc (180:360:0.1);
        \draw (0.1,0) arc (180:360:0.1);
        \draw (0,-0.15) -- (0,-0.25);
        \draw (0,-0.25) arc (180:360:0.05);
        \draw (0,-0.25) arc (0:-180:0.05);
    \end{scope}

    \draw (-1.4,-0.3) .. controls (-2.5,-0.5) and (-2.5,1) .. (-1.5,0.6);
    \draw (-0.5,-0.8) arc (200:340:0.4);
\end{scope}

% --------------------------------------------------
% Layer Symbols
% --------------------------------------------------

\newcommand{\basePos}{2.15,-0.5}

\newcommand{\posZzero}{0.11,2.03}
\newcommand{\scaleZzero}{0.85}
\newcommand{\rotateZzero}{-20}

\newcommand{\posZone}{0.2,2.3}
\newcommand{\scaleZone}{0.8}
\newcommand{\rotateZone}{10}

\newcommand{\posZtwo}{0.32,2.4}
\newcommand{\scaleZtwo}{0.75}
\newcommand{\rotateZtwo}{-40}

\begin{scope}[shift={(\basePos)}]
    \node[draw=none, scale=\scaleZzero, rotate=\rotateZzero]
        at (\posZzero) {$z_0$};

    \begin{scope}[
        shift={(\posZone)},
        scale=\scaleZone,
        rotate=\rotateZone,
        transform shape
    ]
        \node[draw=none] at (0,0) {$z_i$};
    \end{scope}

    \begin{scope}[
        shift={(\posZtwo)},
        scale=\scaleZtwo,
        rotate=\rotateZtwo,
        transform shape
    ]
        \node[draw=none] at (0,0) {$z_5$};
    \end{scope}
\end{scope}

\end{tikzpicture}

\vspace{-4em}
\scriptsize

\noindent
\textbf{Note:}
The state labels and step numbers used in this diagram
correspond to the definitions introduced in
``CASE-A1: Schrödinger's Sisyphean Loop (Rev. 1.6).''
They are reused here for conceptual continuity and do not
represent an exhaustive or formally complete state model.

\section*{Technical Observations}

\begin{itemize}
    \item This diagram is based on the observed behavior of a specific Nikon Z f
    unit (firmware v3.01) and represents a non-exhaustive set of
    user-observable state transitions. Consistent behavior across all units or
    environments is not guaranteed. Actual transitions may depend on previous
    interaction history and cannot be fully determined by visible states alone.

    \item As of firmware version 3.00, Nikon has confirmed that when
    transitioning from S5 to S6 via Step 9, recovery is only possible via the
    DISP button or by performing a reset, and that this behavior is considered
    part of the product specification.

    \item \textbf{Reproduction parameters:}
    For the specific menu settings required to reproduce each transition
    (e.g., Step 9, 28, 37, 45), please refer to
    \href{https://github.com/acta-lgtm/nikon-z-display5-persistence-case-zf/blob/main/docs/DA/Case_A1_Schrodinger.md}
    {\textbf{Case A1: Schrödinger's Sisyphean Loop}}.

    \item \textbf{Indeterminacy of state transitions:}
    The transition from Step 3 is not uniquely determined by the external input
    alone; it is dependent on internal state history that is not externally
    observable. Under these conditions, even if two cameras appear to have
    identical settings and display states, their behavior may diverge due to
    differences in internal state. As a result, when encountering a shooting
    opportunity that requires Live View, it is not possible to determine in
    advance which camera will reliably allow capture.

    For example, in Step 5 the system transitions to S1 (Live View), whereas
    in Step 28 it transitions into the S6 (Info display) loop. Thus, the
    outcome is not uniquely determined for the same input, meaning that whether
    shooting is possible cannot be known prior to operation.

    \item \textbf{Loop behavior and priority:}
    This camera provides several monitor modes, including ``Automatic display
    switching,'' ``Viewfinder only,'' ``Monitor only,'' and ``Prioritize
    viewfinder (1, 2).'' However, once the system enters the dashed loop region
    centered on S6, it cannot return to Live View through the standard user
    input path defined within the normal state transition graph. In this state,
    even if the user has selected ``Prioritize viewfinder,'' the behavior
    effectively becomes equivalent to ``Viewfinder only,'' and Live View on the
    monitor is no longer available.

    \item \textbf{Recovery paths (Step 67):}
    Recovery from the persistent Info-display state was observed through
    DISP-related display cycling while the LCD monitor remained active
    (e.g., S6 with the LCD open). When the LCD was closed (S7), the same input
    did not produce a recovery transition.

    Additional recovery paths were later observed, including:
    (i) reassigning the ``DISP Cycle view info display'' function to a custom
    button and then pressing it, and
    (ii) disabling ``Display 5'' in the custom monitor shooting display
    settings.

    However, under configurations where no DISP-related recovery path remained
    accessible through ordinary shooting operations, the persistent state could
    not be cleared through standard UI interaction alone. Power cycling or
    removing and reinserting the battery was not observed to reliably terminate
    this state.
\end{itemize}

\end{document}
